% The Stata example code and results are generated by running
% "do examples" in Stata.

\author{Christopher James Rose \and Mahima Venkateswaran}
\title{Supplementary Material for:\\ Handling an Unanticipated Intercurrent
Event in a Digital Health Intervention Trial: A Case Study}
\date{}
\maketitle

\section{Introduction}

This document is supplementary material for the accompanying case study of
an unanticipated intercurrent event (ICE) in the eRegMat trial, a
cluster-randomized digital health trial in Bangladesh \citep{rahman2021,
venkateswaran2024}, in which software delays meant some intervention
participants received only part of the intervention. We assume the reader
is familiar with the case study and its estimand framework \citep{iche9r1}.

Section~\ref{sec:methods} describes the simulation used to estimate the
bias, precision, and power of four candidate analysis strategies for the
ICE. Section~\ref{sec:results} reports the results, and
Section~\ref{sec:conclusions} interprets them and states which strategy
was chosen for the eRegMat trial. Section~\ref{sec:application} applies
the chosen strategy to example trial data. Section~\ref{sec:discussion}
discusses alternative analytical approaches.


\section{Methods}\label{sec:methods}

We used simulation to estimate the bias, precision, and power of four
candidate strategies for analyzing the eRegMat trial in the presence of
the ICE. We synthesized 1\,000 realizations of the trial in Stata,
specifying and hence knowing the true treatment effect in each, and
applied each strategy to each realization. Comparing estimates to the
known true effect gives, for each strategy, an estimate of bias
(systematic deviation from the true effect), precision (standard error of
the estimate), and power (the probability of correctly rejecting a false
null hypothesis). We used 1\,000 realizations because standard errors on
simulation estimates (e.g.,~of bias) are proportional to $1/\sqrt{n}$,
where $n$ is the number of realizations: a further $100\times$ increase in
$n$ would improve precision only $10\times$, at $100\times$ the
computational cost, and would be unlikely to change the choice of
strategy.

Stata code is presented below in a literate programming style
\citep{knuth1984}, interleaved with the prose describing it and the output
it produces.

\subsection{Simulating trial data}

We cleared Stata's memory, set a random-number seed for reproducibility,
and set a default numerical display format:

\begin{minipage}{\linewidth}
\begin{stlog}
. clear all
{\smallskip}
. set seed 12345
{\smallskip}
. set cformat \%9.2f
{\smallskip}
\nullskip
\end{stlog}
\end{minipage}

We assumed 59 clinics (clusters), matching the eRegMat trial, with clinic
size drawn from a log-normal distribution (mean~3, SD~1 on the log scale)
chosen to approximate the trial's distribution of clinic sizes:

\begin{minipage}{\linewidth}
\begin{stlog}
. global num_clusters    59
{\smallskip}
. global cluster_sz_mean 3
{\smallskip}
. global cluster_sz_sd   1
{\smallskip}
\nullskip
\end{stlog}
\end{minipage}

Because ``modified Poisson'' regression for risk ratios
operates on the log-risk scale, we specified the control-arm risk and the
treatment effect (log risk ratio) on this scale: a control risk of 15\%
and a treatment effect corresponding to a risk ratio of 0.75, both chosen
independently of the trial data:

\begin{minipage}{\linewidth}
\begin{stlog}
. global log_control_risk = log(0.15)
{\smallskip}
. global log_rr           = log(0.75)
{\smallskip}
\nullskip
\end{stlog}
\end{minipage}

To reflect the cluster-randomized design, we assumed a between-clinic
variance on the log-risk scale, based on comparable trials:

\begin{minipage}{\linewidth}
\begin{stlog}
. global cluster_var 0.07
{\smallskip}
\nullskip
\end{stlog}
\end{minipage}

The intervention was affected by the unanticipated ICE for some
participants; we also call this \emph{disruption}, following the variable
name used for it (Section~\ref{sec:application}). Because participants
were enrolled on a rolling basis and the missing software components
became available partway through recruitment, disruption depended on each
participant's enrollment timing rather than their clinic: intervention
participants enrolled before the components became available were
disrupted, and those enrolled after were not. We therefore simulated
disruption as an independent event for each intervention participant
rather than as a clinic-level attribute. We assumed that 63\% of
intervention participants were affected, the proportion recorded during
trial monitoring; because the monitoring was of treatment delivery rather
than outcome, it could inform the simulation without accessing outcome
data or risking unblinding:

\begin{minipage}{\linewidth}
\begin{stlog}
. global nondisrupted_prop  = 0.37
{\smallskip}
\nullskip
\end{stlog}
\end{minipage}

Two Stata programs generate one realization of the trial.
\stcmd{simulate\_clusters} creates one row per clinic, with a random
cluster effect, a random size, and 1:1 randomization to intervention or
control:

\begin{minipage}{\linewidth}
\begin{stlog}
. program simulate_clusters, nclass
  1.   version 18
  2.   quietly set obs \$num_clusters
  3.   generate byte cluster     = _n
  4.   generate log_cluster_risk = rnormal(0, sqrt(\$cluster_var))
  5.   generate log_cluster_size = rnormal(\$cluster_sz_mean, \$cluster_sz_sd)
  6.   splitsample , generate(intervention) values(0 1) rround
  7.   label variable cluster          "Cluster identifier (observable)"
  8.   label variable log_cluster_risk "Log cluster risk (unobservable)"
  9.   label variable log_cluster_size "Log cluster size (observable)"
 10.   label variable intervention     "Intervention (observable)"
 11.   label define   intervention     0 "Control"                          ///
>                                      1 "Intervention"
 12.   label values   intervention     intervention
 13. end
{\smallskip}
\nullskip
\end{stlog}
\end{minipage}

\stcmd{simulate\_participants} expands the clinic-level data to one row
per participant and simulates each participant's outcome as a Bernoulli
draw whose probability is given by the control risk, their clinic's
random effect, and, for intervention participants not affected by the
ICE, the treatment effect:

\begin{minipage}{\linewidth}
\begin{stlog}
. program simulate_participants, nclass
  1.   version 18
  2.   quietly expand exp(log_cluster_size)
  3.   generate nondisrupted = intervention *                               ///
>                              (runiform() < \$nondisrupted_prop)
  4.   generate log_risk     = \$log_control_risk + log_cluster_risk +       ///
>                              (nondisrupted * \$log_rr)
  5.   generate outcome      = rbinomial(1, exp(log_risk))
  6.   label variable nondisrupted "Not disrupted (observable)"
  7.   label define   nondisrupted 0 "Disrupted"                            ///
>                                  1 "Not disrupted"
  8.   label values   nondisrupted nondisrupted
  9.   label variable log_risk     "Log risk (unobservable)"
 10.   label variable outcome      "Outcome (observable)"
 11.   label define   outcome      0 "No morbidity or mortality"            ///
>                                  1 "Morbidity or mortality"
 12.   label values   outcome      outcome
 13. end
{\smallskip}
\nullskip
\end{stlog}
\end{minipage}

\subsection{Strategy 1: Treatment policy}

Each strategy is a modified Poisson regression for the log risk of the
outcome (severe morbidity or mortality), $Y_i$, for participant $i$, with
a random intercept $\gamma_i$ for participant $i$'s clinic, to account for
the cluster-randomized design. Strategy 1 corresponds to the eRegMat
trial's original analysis, planned before any ICE-handling strategy had
been specified; in ICH~E9(R1) terms, it is the treatment policy strategy,
analyzing all randomized participants according to their allocated
treatment (intention-to-treat, ITT) without adjusting for the ICE. It is
implemented by equation~\eqref{strategy1}:
%
\begin{equation}\label{strategy1}
  \log \mathrm{E}(Y_i) = \alpha + \beta x_i + \gamma_i \quad \mathrm{where}\quad i \in \mathrm{ITT}
\end{equation}
%
where $x_i$ indicates participant $i$'s randomized treatment
($x_i\!=\!0\!\Rightarrow\!\mathrm{control}$,
$x_i\!=\!1\!\Rightarrow\!\mathrm{intervention}$) and $\beta$, the treatment
effect (log risk ratio), is the quantity of interest. The program
\stcmd{prespecified\_analysis}, so named because it was the trial's
original analysis, returns the point estimate, standard error, and a test
of no treatment effect (at the 5\% level) for use in the simulation:

\begin{minipage}{\linewidth}
\begin{stlog}
. program prespecified_analysis, rclass
  1.   version 18
  2.   mepoisson outcome i.intervention ||cluster:, vce(robust)
  3.   return scalar es1_Prespecified      = _b[1.intervention]
  4.   return scalar se1_Prespecified      = _se[1.intervention]
  5.   return scalar reject1_Prespecified  =                                ///
>                     r(table)["pvalue", "outcome:1.intervention"] < 0.05
  6. end
{\smallskip}
\nullskip
\end{stlog}
\end{minipage}

\subsection{Strategy 2: Modified intention-to-treat}

Strategy 2 restricts the analysis set to control participants and
intervention participants not affected by the ICE; the model is otherwise
identical to Strategy 1, giving equation~\eqref{strategy2}:
%
\begin{equation}\label{strategy2}
  \log \mathrm{E}(Y_i) = \alpha + \beta x_i + \gamma_i \quad \mathrm{where}\quad i \in \mathrm{mITT}
\end{equation}
%
where $\mathrm{mITT}$ denotes the restricted analysis set. The program
\stcmd{mITT\_analysis} restricts \stcmd{mepoisson}'s estimation sample
and otherwise returns the same quantities as \stcmd{prespecified\_analysis}:

\begin{minipage}{\linewidth}
\begin{stlog}
. program mITT_analysis, rclass
  1.   version 18
  2.   local included !intervention | nondisrupted
  3.   mepoisson outcome i.intervention if `included' ||cluster: , vce(robust)
  4.   return scalar es2_mITT      = _b[1.intervention]
  5.   return scalar se2_mITT      = _se[1.intervention]
  6.   return scalar reject2_mITT  =                                        ///
>                     r(table)["pvalue", "outcome:1.intervention"] < 0.05
  7. end
{\smallskip}
\nullskip
\end{stlog}
\end{minipage}

\subsection{Strategy 3: As-treated}

Strategy 3 uses the same analysis set as Strategy 1 but classifies
participants by the treatment they actually received rather than the
treatment to which they were randomized, breaking randomization; this
gives equation~\eqref{strategy3}:
%
\begin{equation}\label{strategy3}
  \log \mathrm{E}(Y_i) = \alpha + \beta z_i + \gamma_i \quad \mathrm{where}\quad i \in \mathrm{ITT}
\end{equation}
%
where $z_i$ indicates whether participant $i$ actually received the
intervention ($z_i\!=\!0\!\Rightarrow\!\mathrm{not~received}$,
$z_i\!=\!1\!\Rightarrow\!\mathrm{received}$). The program
\stcmd{astreated\_analysis} estimates a coefficient for
\stcmd{nondisrupted} (which, by construction, indicates whether the
intervention was actually received) rather than for \stcmd{intervention}:

\begin{minipage}{\linewidth}
\begin{stlog}
. program astreated_analysis, rclass
  1.   version 18
  2.   mepoisson outcome i.nondisrupted ||cluster:, vce(robust)
  3.   return scalar es3_AsTreated      = _b[1.nondisrupted]
  4.   return scalar se3_AsTreated      = _se[1.nondisrupted]
  5.   return scalar reject3_AsTreated  =                                 ///
>                       r(table)["pvalue", "outcome:1.nondisrupted"] < 0.05
  6. end
{\smallskip}
\nullskip
\end{stlog}
\end{minipage}

\subsection{Strategy 4: Hypothetical}\label{sec:strategy4}

Strategy 4 estimates the treatment effect under the counterfactual
scenario of full intervention delivery --- the quantity the trial was
designed to answer --- applying the ICH~E9(R1) hypothetical strategy
\citep{iche9r1}. The model decomposes the log risk ratio into the effect
of randomization to intervention ($\beta$) and the effect of receiving
the intervention in full ($\delta$); the originally specified treatment
effect is $\beta+\delta$. Unlike Strategy 3, this strategy uses the ITT
analysis set and preserves randomization, giving equation~\eqref{strategy4}:
%
\begin{equation}\label{strategy4}
  \log \mathrm{E}(Y_i) = \alpha + \beta x_i + \delta x_i^\ast + \gamma_i \quad \mathrm{where}\quad i \in \mathrm{ITT}
\end{equation}
%
where $x_i$ is defined as for Strategy 1 and $x_i^\ast$ indicates whether
participant $i$ received the intervention in full
($x_i^\ast\!=\!0\!\Rightarrow\!\mathrm{partial}$,
$x_i^\ast\!=\!1\!\Rightarrow\!\mathrm{full}$). Because $\beta$ is the
effect of randomization to intervention among participants for whom it
was disrupted, and $\delta$ is the additional effect of receiving it in
full, their sum $\beta+\delta$ --- not $\beta$ or $\delta$ alone ---
estimates the effect of full intervention delivery relative to control,
the quantity of interest. The program \stcmd{hypothetical\_analysis} uses
\stcmd{lincom} to estimate $\beta+\delta$ with its associated standard
error and hypothesis test:

\begin{minipage}{\linewidth}
\begin{stlog}
. program hypothetical_analysis, rclass
  1.   version 18
  2.   mepoisson outcome i.intervention i.nondisrupted ||cluster:, vce(robust)
  3.   lincom _b[1.intervention] + _b[1.nondisrupted]
  4.   return scalar es4_Hypothetical      = r(estimate)
  5.   return scalar se4_Hypothetical      = r(se)
  6.   return scalar reject4_Hypothetical  = r(p) < 0.05
  7. end
{\smallskip}
\nullskip
\end{stlog}
\end{minipage}

\subsection{Running the simulation}

The program \stcmd{simulate\_trial} composes the programs above into a
single realization of the trial, applying all four strategies:

\begin{minipage}{\linewidth}
\begin{stlog}
. program simulate_trial, rclass
  1.   version 18
  2.   clear
  3.   simulate_clusters
  4.   simulate_participants
  5.   foreach method in prespecified mITT astreated hypothetical {\lbr}
  6.     `method'_analysis
  7.     return add
  8.   {\rbr}
  9.   ereturn clear
 10. end
{\smallskip}
\nullskip
\end{stlog}
\end{minipage}

We simulated 1\,000 realizations:

\begin{minipage}{\linewidth}
\begin{stlog}
. simulate , reps(1000) nolegend : simulate_trial
{\smallskip}
Simulations (1,000): .........10.........20.........30.........40.........50....
> .....60.........70.........80.........90.........100.........110.........120..
> .......130.........140.........150.........160.........170.........180........
> .190.....{\bftt{x}}...200.........210.........220.........230....{\bftt{x}}....240.........250..
> .......260........{\bftt{x}}270.........280.........290.........300.........310........
> .320.........330.........340.........350.........360.........370.........380..
> .......390.........400.........410.........420.........430.........440........
> .450.........460.........470.........480.........490.........500.........510..
> .......520.........530.........540.........550.........560.........570........
> .580.........590.........600.........610.........620.........630.........640..
> .......650.........660.........670.........680.........690.........700.......{\bftt{x}}
> .710.........720.........730.........740.........750.......{\bftt{x}}.760.........770..
> .......780.........790.........800.........810.........820.........830........
> .840.........850.........860.........870.........880.........890.........900..
> .......910.........920.........930.........940.........950.........960........
> .{\bftt{x}}.........980.........990.........1,000 done
{\bftt{x}}: Error occurred when {\bftt{simulate}} executed {\bftt{simulate_trial}}.
\nullskip
\end{stlog}
\end{minipage}

All four strategies converged for the great majority of realizations; we
did not investigate the small number of convergence failures further.


\section{Results}\label{sec:results}

We reshaped the simulation output into long format, creating a
\stcmd{method} factor variable, to estimate bias, precision, and power for
each strategy:

\begin{minipage}{\linewidth}
\begin{stlog}
. generate simulation = _n
{\smallskip}
. tempvar method
{\smallskip}
. reshape long es se reject , i(simulation) j(`method') string
(j = 1_Prespecified 2_mITT 3_AsTreated 4_Hypothetical)
{\smallskip}
Data                               Wide   ->   Long
\HLI{77}
Number of observations            1,000   ->   4,000       
Number of variables                  13   ->   5           
j variable (4 values)                     ->   __000000
xij variables:
es1_Prespecified es2_mITT ... es4_Hypothetical-> es
se1_Prespecified se2_mITT ... se4_Hypothetical-> se
reject1_Prespecified reject2_mITT ... reject4_Hypothetical->reject
\HLI{77}
{\smallskip}
. encode `method', generate(method)
{\smallskip}
\nullskip
\end{stlog}
\end{minipage}

\subsection{Bias}

We calculated bias as the percentage difference between each strategy's
estimate and the assumed treatment effect, on the log-risk scale, then
estimated mean bias for each strategy by regression:

\begin{minipage}{\linewidth}
\begin{stlog}
. generate bias = 100 * ((\$log_rr - es) / \$log_rr)
(24 missing values generated)
{\smallskip}
. bysort method: regress bias , noheader
{\smallskip}
\HLI{80}
-> method = 1_Prespecified
\HLI{13}{\TOPT}\HLI{64}
        bias {\VBAR} Coefficient  Std. err.      t    P>|t|     [95\% conf. interval]
\HLI{13}{\PLUS}\HLI{64}
       _cons {\VBAR}      65.00       1.61    40.48   0.000        61.85       68.15
\HLI{13}{\BOTT}\HLI{64}
{\smallskip}
\HLI{80}
-> method = 2_mITT
\HLI{13}{\TOPT}\HLI{64}
        bias {\VBAR} Coefficient  Std. err.      t    P>|t|     [95\% conf. interval]
\HLI{13}{\PLUS}\HLI{64}
       _cons {\VBAR}      -3.61       2.18    -1.66   0.097        -7.88        0.66
\HLI{13}{\BOTT}\HLI{64}
{\smallskip}
\HLI{80}
-> method = 3_AsTreated
\HLI{13}{\TOPT}\HLI{64}
        bias {\VBAR} Coefficient  Std. err.      t    P>|t|     [95\% conf. interval]
\HLI{13}{\PLUS}\HLI{64}
       _cons {\VBAR}      -3.59       1.92    -1.87   0.061        -7.35        0.17
\HLI{13}{\BOTT}\HLI{64}
{\smallskip}
\HLI{80}
-> method = 4_Hypothetical
\HLI{13}{\TOPT}\HLI{64}
        bias {\VBAR} Coefficient  Std. err.      t    P>|t|     [95\% conf. interval]
\HLI{13}{\PLUS}\HLI{64}
       _cons {\VBAR}      -3.70       2.15    -1.72   0.086        -7.91        0.52
\HLI{13}{\BOTT}\HLI{64}
\nullskip
\end{stlog}
\end{minipage}

Strategy 1 (Treatment policy) was heavily biased: 65\% ($P<0.001$) towards
$+\infty$ on the log-risk scale, corresponding to an average estimated
risk ratio of 0.90 against an assumed true value of 0.75. Because risk
ratios below one favor the intervention, Strategy 1 would be expected to
substantially underestimate the intervention's benefit, or to suggest harm
where none exists. Strategies 2--4 showed no evidence of bias ($P>0.05$ in
all cases).

\subsection{Precision}

We estimated mean standard error (precision) using the same approach, to
four significant figures given the small magnitudes involved:

\begin{minipage}{\linewidth}
\begin{stlog}
. bysort method: regress se , cformat(\%9.4g) noheader
{\smallskip}
\HLI{80}
-> method = 1_Prespecified
\HLI{13}{\TOPT}\HLI{64}
          se {\VBAR} Coefficient  Std. err.      t    P>|t|     [95\% conf. interval]
\HLI{13}{\PLUS}\HLI{64}
       _cons {\VBAR}      .1412   .0006668   211.74   0.000        .1399       .1425
\HLI{13}{\BOTT}\HLI{64}
{\smallskip}
\HLI{80}
-> method = 2_mITT
\HLI{13}{\TOPT}\HLI{64}
          se {\VBAR} Coefficient  Std. err.      t    P>|t|     [95\% conf. interval]
\HLI{13}{\PLUS}\HLI{64}
       _cons {\VBAR}       .187    .001059   176.65   0.000        .1849       .1891
\HLI{13}{\BOTT}\HLI{64}
{\smallskip}
\HLI{80}
-> method = 3_AsTreated
\HLI{13}{\TOPT}\HLI{64}
          se {\VBAR} Coefficient  Std. err.      t    P>|t|     [95\% conf. interval]
\HLI{13}{\PLUS}\HLI{64}
       _cons {\VBAR}      .1607    .001086   147.99   0.000        .1586       .1628
\HLI{13}{\BOTT}\HLI{64}
{\smallskip}
\HLI{80}
-> method = 4_Hypothetical
\HLI{13}{\TOPT}\HLI{64}
          se {\VBAR} Coefficient  Std. err.      t    P>|t|     [95\% conf. interval]
\HLI{13}{\PLUS}\HLI{64}
       _cons {\VBAR}      .1859    .001002   185.53   0.000         .184       .1879
\HLI{13}{\BOTT}\HLI{64}
\nullskip
\end{stlog}
\end{minipage}

Strategy 1 had the smallest (best) standard error, despite being the most
biased, because it does not model the additional uncertainty introduced
by the ICE (i.e.,~what outcomes would have occurred had it not occurred).
Strategy 3 (As-treated) was more precise than Strategies 2 and 4, between
which there was no evidence of a difference in precision:

\begin{minipage}{\linewidth}
\begin{stlog}
. regress se i.method , nocons noheader
\HLI{14}{\TOPT}\HLI{64}
           se {\VBAR} Coefficient  Std. err.      t    P>|t|     [95\% conf. interval]
\HLI{14}{\PLUS}\HLI{64}
       method {\VBAR}
      2_mITT  {\VBAR}       0.19       0.00    76.64   0.000         0.18        0.19
 3_AsTreated  {\VBAR}       0.16       0.00    65.85   0.000         0.16        0.17
4_Hypothet{\tytilde}l  {\VBAR}       0.19       0.00    76.20   0.000         0.18        0.19
\HLI{14}{\BOTT}\HLI{64}
{\smallskip}
. test 2.method = 4.method
{\smallskip}
 ( 1)  2.method - 4.method = 0
{\smallskip}
       F(  1,  3973) =    0.10
            Prob > F =    0.7552
{\smallskip}
\nullskip
\end{stlog}
\end{minipage}

\subsection{Power}

We estimated power, the probability of correctly rejecting a false null
hypothesis, using exact (Clopper--Pearson) confidence intervals for each
strategy's rejection proportion:

\begin{minipage}{\linewidth}
\begin{stlog}
. bysort method: ci proportions reject
{\smallskip}
\HLI{80}
-> method = 1_Prespecified
{\smallskip}
                                                            Binomial exact   
    Variable {\VBAR}        Obs  Proportion    Std. err.       [95\% conf. interval]
\HLI{13}{\PLUS}\HLI{63}
      reject {\VBAR}        994    .1197183    .0102967        .1001838    .1415362
{\smallskip}
\HLI{80}
-> method = 2_mITT
{\smallskip}
                                                            Binomial exact   
    Variable {\VBAR}        Obs  Proportion    Std. err.       [95\% conf. interval]
\HLI{13}{\PLUS}\HLI{63}
      reject {\VBAR}        994    .3420523     .015047        .3125625    .3724812
{\smallskip}
\HLI{80}
-> method = 3_AsTreated
{\smallskip}
                                                            Binomial exact   
    Variable {\VBAR}        Obs  Proportion    Std. err.       [95\% conf. interval]
\HLI{13}{\PLUS}\HLI{63}
      reject {\VBAR}        994    .4567404    .0157996        .4254383    .4882995
{\smallskip}
\HLI{80}
-> method = 4_Hypothetical
{\smallskip}
                                                            Binomial exact   
    Variable {\VBAR}        Obs  Proportion    Std. err.       [95\% conf. interval]
\HLI{13}{\PLUS}\HLI{63}
      reject {\VBAR}        994    .3490946    .0151195         .319442    .3796442
\nullskip
\end{stlog}
\end{minipage}

Strategy 1's power was at most 14\% (upper confidence bound). Strategies
2--4 were substantially more powerful, but still only modestly powered, due to
the substantial proportion of participants affected by the ICE. Strategy 3
(As-treated) was the most powerful; there was no evidence of a difference
between Strategies 2 and 4:

\begin{minipage}{\linewidth}
\begin{stlog}
. regress reject i.method , nocons noheader
\HLI{14}{\TOPT}\HLI{64}
       reject {\VBAR} Coefficient  Std. err.      t    P>|t|     [95\% conf. interval]
\HLI{14}{\PLUS}\HLI{64}
       method {\VBAR}
      2_mITT  {\VBAR}       0.34       0.01    23.81   0.000         0.31        0.37
 3_AsTreated  {\VBAR}       0.46       0.01    31.79   0.000         0.43        0.48
4_Hypothet{\tytilde}l  {\VBAR}       0.35       0.01    24.30   0.000         0.32        0.38
\HLI{14}{\BOTT}\HLI{64}
{\smallskip}
. test 2.method = 4.method
{\smallskip}
 ( 1)  2.method - 4.method = 0
{\smallskip}
       F(  1,  3973) =    0.12
            Prob > F =    0.7289
{\smallskip}
\nullskip
\end{stlog}
\end{minipage}


\section{Conclusions}\label{sec:conclusions}

Strategy 1 (Treatment policy) was unsuitable: it was heavily biased and,
because it does not reflect the uncertainty introduced by the ICE, gave a
misleadingly precise answer to a question the trial was not designed to
answer. Strategies 2--4 showed no evidence of bias. Strategy 3
(As-treated) was the most powerful of these, but breaks randomization,
reintroducing the risk of confounding that randomization is designed to
eliminate, so its estimate cannot be given a strong causal interpretation.
Strategies 2 (Modified intention-to-treat) and 4 (Hypothetical) had
similar power, with no evidence of a difference between them.

We selected Strategy 4 for the eRegMat trial. Like Strategy 2, it
preserves randomization and the intention-to-treat principle; unlike
Strategy 2, its standard error reflects the uncertainty introduced by the
ICE, and it estimates the treatment effect originally specified by the
trial under the counterfactual scenario of full intervention delivery,
rather than an effect restricted to a subset of participants. This
choice, and the simulation results supporting it, were finalized before
the trial data were analyzed or unblinded.


\section{Application to example trial data}\label{sec:application}

The accompanying file \stcmd{disrupted.dta} contains fictitious data with
the same structure and variable names used throughout this document: one
row per participant, with variables \stcmd{cluster}, \stcmd{intervention},
\stcmd{nondisrupted}, and \stcmd{outcome}.

\begin{minipage}{\linewidth}
\begin{stlog}
. use disrupted , clear
{\smallskip}
. codebook
{\smallskip}
\HLI{80}
cluster                                          Cluster identifier (observable)
\HLI{80}
{\smallskip}
                  Type: Numeric (byte)
{\smallskip}
                 Range: [1,59]                        Units: 1
         Unique values: 59                        Missing .: 0/2,356
{\smallskip}
                  Mean: 33.1668
             Std. dev.: 17.5358
{\smallskip}
           Percentiles:     10\%       25\%       50\%       75\%       90\%
                             10        19        33        53        53
{\smallskip}
\HLI{80}
intervention                                           Intervention (observable)
\HLI{80}
{\smallskip}
                  Type: Numeric (byte)
                 Label: intervention
{\smallskip}
                 Range: [0,1]                         Units: 1
         Unique values: 2                         Missing .: 0/2,356
{\smallskip}
            Tabulation: Freq.   Numeric  Label
                          991         0  Control
                        1,365         1  Intervention
{\smallskip}
\HLI{80}
nondisrupted                                          Not disrupted (observable)
\HLI{80}
{\smallskip}
                  Type: Numeric (float)
                 Label: nondisrupted
{\smallskip}
                 Range: [0,1]                         Units: 1
         Unique values: 2                         Missing .: 0/2,356
{\smallskip}
            Tabulation: Freq.   Numeric  Label
                        1,844         0  Disrupted
                          512         1  Not disrupted
{\smallskip}
\HLI{80}
outcome                                                     Outcome (observable)
\HLI{80}
{\smallskip}
                  Type: Numeric (float)
                 Label: outcome
{\smallskip}
                 Range: [0,1]                         Units: 1
         Unique values: 2                         Missing .: 0/2,356
{\smallskip}
            Tabulation: Freq.   Numeric  Label
                        2,011         0  No morbidity or mortality
                          345         1  Morbidity or mortality
\nullskip
\end{stlog}
\end{minipage}

Clinics, not participants, were randomized, so the arms are not exactly
balanced. Of participants randomized to intervention, 62\% were affected
by the ICE; the apparent disruption of control participants below is an
artifact of how control participants are coded, since only the
intervention could be disrupted:

\begin{minipage}{\linewidth}
\begin{stlog}
. tabulate intervention nondisrupted , row
{\smallskip}
{\TLC}\HLI{16}{\TRC}
{\VBAR} Key            {\VBAR}
{\LFTT}\HLI{16}{\RGTT}
{\VBAR}   {\sltt{frequency}}    {\VBAR}
{\VBAR} {\sltt{row percentage}} {\VBAR}
{\BLC}\HLI{16}{\BRC}
{\smallskip}
             {\VBAR}     Not disrupted
Intervention {\VBAR}     (observable)
(observable) {\VBAR} Disrupted  Not disru {\VBAR}     Total
\HLI{13}{\PLUS}\HLI{22}{\PLUS}\HLI{10}
     Control {\VBAR}       991          0 {\VBAR}       991 
             {\VBAR}    100.00       0.00 {\VBAR}    100.00 
\HLI{13}{\PLUS}\HLI{22}{\PLUS}\HLI{10}
Intervention {\VBAR}       853        512 {\VBAR}     1,365 
             {\VBAR}     62.49      37.51 {\VBAR}    100.00 
\HLI{13}{\PLUS}\HLI{22}{\PLUS}\HLI{10}
       Total {\VBAR}     1,844        512 {\VBAR}     2,356 
             {\VBAR}     78.27      21.73 {\VBAR}    100.00 
{\smallskip}
\nullskip
\end{stlog}
\end{minipage}

For comparison, we also apply Strategy 1 (Treatment policy) --- the
trial's original, prespecified analysis --- to these data. This
comparison was not used to choose an analysis strategy: the choice of
Strategy 4 (Section~\ref{sec:conclusions}) was finalized by simulation
alone, before the real eRegMat trial data were analyzed or unblinded.

\begin{minipage}{\linewidth}
\begin{stlog}
. mepoisson outcome i.intervention ||cluster:, vce(robust) eform(RR) nolog
{\smallskip}
Mixed-effects Poisson regression                Number of obs     =      2,356
Group variable: cluster                         Number of groups  =         59
{\smallskip}
                                                Obs per group:
                                                              min =          2
                                                              avg =       39.9
                                                              max =        412
{\smallskip}
Integration method: mvaghermite                 Integration pts.  =          7
{\smallskip}
                                                Wald chi2(1)      =       0.01
Log pseudolikelihood = -1007.3197               Prob > chi2       =     0.9214
                                (Std. err. adjusted for 59 clusters in cluster)
\HLI{14}{\TOPT}\HLI{64}
              {\VBAR}               Robust
      outcome {\VBAR}         RR   std. err.      z    P>|z|     [95\% conf. interval]
\HLI{14}{\PLUS}\HLI{64}
 intervention {\VBAR}
Intervention  {\VBAR}       1.01       0.15     0.10   0.921         0.76        1.35
        _cons {\VBAR}       0.15       0.02   -18.19   0.000         0.12        0.18
\HLI{14}{\PLUS}\HLI{64}
cluster       {\VBAR}
    var(_cons){\VBAR}       0.03       0.04                          0.00        0.39
\HLI{14}{\BOTT}\HLI{64}
Note: Estimates are transformed only in the first equation.
\nullskip
\end{stlog}
\end{minipage}

Strategy 1 estimates the risk ratio to be 1.01 (95\%~CI~0.76 to 1.35;
$P=0.921$). As anticipated from the simulation results, this estimate is
likely biased and its apparent precision misleading.

We now apply the chosen strategy, Strategy 4 (Section~\ref{sec:strategy4}):

\begin{minipage}{\linewidth}
\begin{stlog}
. use disrupted, clear
{\smallskip}
. mepoisson outcome i.intervention i.nondisrupted ||cluster:, vce(robust) nolog
{\smallskip}
Mixed-effects Poisson regression                Number of obs     =      2,356
Group variable: cluster                         Number of groups  =         59
{\smallskip}
                                                Obs per group:
                                                              min =          2
                                                              avg =       39.9
                                                              max =        412
{\smallskip}
Integration method: mvaghermite                 Integration pts.  =          7
{\smallskip}
                                                Wald chi2(2)      =       0.74
Log pseudolikelihood = -1007.0253               Prob > chi2       =     0.6918
                                (Std. err. adjusted for 59 clusters in cluster)
\HLI{14}{\TOPT}\HLI{64}
              {\VBAR}               Robust
      outcome {\VBAR} Coefficient  std. err.      z    P>|z|     [95\% conf. interval]
\HLI{14}{\PLUS}\HLI{64}
 intervention {\VBAR}
Intervention  {\VBAR}       0.06       0.16     0.37   0.712        -0.25        0.37
              {\VBAR}
 nondisrupted {\VBAR}
Not disrup..  {\VBAR}      -0.11       0.13    -0.86   0.391        -0.38        0.15
        _cons {\VBAR}      -1.90       0.10   -18.19   0.000        -2.11       -1.70
\HLI{14}{\PLUS}\HLI{64}
cluster       {\VBAR}
    var(_cons){\VBAR}       0.03       0.04                          0.00        0.36
\HLI{14}{\BOTT}\HLI{64}
\nullskip
\end{stlog}
%
\vspace{1em}
\begin{stlog}
. lincom _b[1.intervention] + _b[1.nondisrupted] , eform
{\smallskip}
 ( 1)  [outcome]1.intervention + [outcome]1.nondisrupted = 0
{\smallskip}
\HLI{13}{\TOPT}\HLI{64}
     outcome {\VBAR}     exp(b)   Std. err.      z    P>|z|     [95\% conf. interval]
\HLI{13}{\PLUS}\HLI{64}
         (1) {\VBAR}       0.95       0.15    -0.35   0.729         0.69        1.30
\HLI{13}{\BOTT}\HLI{64}
\nullskip
\end{stlog}
\end{minipage}

Strategy 4 estimates the risk ratio to be 0.95 (95\%~CI~0.69 to 1.30;
$P=0.729$). Unlike Strategy 1, the point estimate suggests the
intervention is beneficial, consistent with the direction of bias the
simulation identified for Strategy 1. The width of this confidence
interval is consistent with the modest power estimated for Strategy 4
(Section~\ref{sec:results}).


\section{Discussion}\label{sec:discussion}

Other approaches exist for estimating treatment effects when some
participants do not receive their allocated treatment. The Heckman
selection model \citep{heckman1979} corrects for non-random treatment
allocation by modelling each participant's probability of receiving the
treatment they received. We did not use it because clinics, and hence
participants, were allocated to the disrupted or undisrupted intervention
entirely at random, so there is no selection process to model. Local
average treatment effect (LATE, also called complier average causal
effect) methods \citep{angrist1995,troncoso2022} estimate the effect
among participants who adhere to treatment, a different quantity to the
ITT effect estimated by Strategy 4. The eRegMat trial prespecified an ITT
estimand; LATE would be relevant were the research question to concern
the effect of adherence itself.


\bibliographystyle{firstauthor}
\bibliography{bibliography}

\endinput
